\chapter{The Hybrid Type System}

The hybrid system in \repo{src/deppy/hybrid/} is best understood as a lift of the plain
system into a layered, trust-bearing semantics. The key file is
\repo{src/deppy/hybrid/core/types.py}, whose opening documentation states the central
mathematical object directly:
\[
  \HybridTy : (\Site(P) \times \Layer)^{\op} \to \mathsf{RefinementLattice} \times \mathsf{TrustLattice}.
\]
This formula does a large amount of conceptual work. It says that a hybrid type is not
just a base type plus an extra annotation. It is a section defined over a product site
whose second coordinate ranges over the ordered set
\[
  \mathrm{INTENT} < \mathrm{STRUCTURAL} < \mathrm{SEMANTIC} < \mathrm{PROOF} < \mathrm{CODE}.
\]

\section*{Layered sections}
The hybrid core defines separate section-like structures for each layer. Intent sections
store natural-language descriptions and structured claims. Structural sections store
machine-checkable predicates, often tied to Z3 or runtime evaluation. Semantic sections
store oracle-judged predicates with prompts, rubrics, and confidence thresholds. Proof
sections store Lean-facing theorem content. Code sections store executable source and
structural program representations.

The result is a model in which one and the same semantic object can be described at
multiple levels without collapsing those levels into one another. This is the central
advantage of the hybrid formulation. It gives \DepPy{} a disciplined answer to the
problem of mixed evidence.

\begin{figure}[h]
\centering
\begin{tikzpicture}[
  layer/.style={draw=chapterblue, rounded corners=2pt, fill=chapterblue!6, minimum width=6.2cm, minimum height=0.8cm, align=center},
  score/.style={draw=accentgold, rounded corners=2pt, fill=accentgold!12, minimum width=3.4cm, minimum height=0.8cm, align=center},
  >=Latex
]
\node[layer] (intent) {INTENT \\ natural-language thesis and desired behavior};
\node[layer, below=0.45cm of intent] (structural) {STRUCTURAL \\ decidable predicates, Z3 obligations, runtime checks};
\node[layer, below=0.45cm of structural] (semantic) {SEMANTIC \\ oracle-judged rubrics, confidence, provenance};
\node[layer, below=0.45cm of semantic] (proof) {PROOF \\ Lean translations, theorem objects, residual assumptions};
\node[layer, below=0.45cm of proof] (code) {CODE \\ executable Python, extracted artifacts, assembled modules};
\node[score, right=1.7cm of semantic] (trust) {trust lattice \\ weakest-link composition};
\foreach \a/\b in {intent/structural,structural/semantic,semantic/proof,proof/code} {
  \draw[->, thick, chapterblue] (\a) -- (\b);
}
\draw[->, thick, accentgold] (semantic) -- (trust);
\draw[->, thick, accentgold] (proof.east) .. controls +(1.1,0.0) and +(0.0,-0.8) .. (trust.south);
\end{tikzpicture}
\caption{The layered object behind hybrid typing in \DepPy{}. The system keeps intent, structure, semantics, proof, and code distinct while forcing them to cohere under restriction and trust composition.}
\end{figure}

\begin{definition}[Hybrid type]
A hybrid type is a tuple of layer sections
\[
  \tau = (\tau_I, \tau_S, \tau_M, \tau_P, \tau_C)
\]
with components for intent, structural, semantic, proof, and code layers, together with
restriction maps between adjacent layers and a trust label that records the strongest
validated interpretation currently available.
\end{definition}

\section*{Dependent formers survive the lift}
The hybrid layer does not abandon the ordinary dependent type formers. Instead it defines
hybrid analogues of them. The repository names \py{HybridPiType}, \py{HybridSigmaType},
and \py{HybridRefinementType}. Their point is not to replace $\Pi$ and $\Sigma$ with new
syntax, but to reinterpret them in a setting where domain and codomain can themselves carry
structural and semantic predicates, provenance, and trust labels.

This is why the hybrid system should be read as an extension rather than a fork. A hybrid
$\Pi$-type is still fundamentally a dependent function. A hybrid $\Sigma$-type is still a
witness-bearing or decomposition-bearing pair. A hybrid refinement type still says
``values of this type satisfy this predicate,'' except that the predicate may now split into
a decidable part and a semantic part:
\[
  \{x : T \mid \varphi_d(x) \land \varphi_s(x)\}.
\]

\section*{Coherence as vanishing first cohomology}
The hybrid core says a type is coherent when first cohomology vanishes:
\[
  \Hc{1}(\Layer, \HybridTy) = 0.
\]
Operationally, this means adjacent layers agree on their overlaps. A function whose code
runs but whose semantic layer rejects its behavior is not coherent. A proof artifact that
postulates a theorem stronger than what the structural layer can justify is not coherent.
A natural-language intent that is never realized by any code or proof is not coherent.

This criterion is philosophically valuable because it makes mismatch visible without
pretending every mismatch is the same kind of failure. Some failures are structural. Some
are semantic. Some are proof-level. The language of layered cocycles keeps those failures
organized.

The hybrid geometry can also be pictured as a commutative square: restrictions in the site
dimension should commute with forgetful transport between layers. For a local chart overlap
$U \subseteq V$, one wants a square of the form
\[
\begin{tikzcd}[column sep=large, row sep=large]
\tau_S(V) \arrow[r, "\rho_{I \leftarrow S}"] \arrow[d, "\mathrm{res}_{U,V}"'] &
\tau_I(V) \arrow[d, "\mathrm{res}_{U,V}"] \\
\tau_S(U) \arrow[r, "\rho_{I \leftarrow S}"] &
\tau_I(U)
\end{tikzcd}
\]
to commute. When such squares fail to commute across a cover, the mismatch is precisely
the sort of obstruction that the hybrid cohomology language is meant to expose.

\begin{remark}
The repository sometimes introduces convenience trust enums in one file and richer trust
structures in another. A precise reading is that the hybrid package exposes a lightweight
public face and a heavier core. The paper should acknowledge this rather than flatten it.
\end{remark}
